\subsection{梯度裁剪}

梯度裁剪（gradient clipping）是防止训练不稳定的重要工具。最常见的做法是\textbf{全局梯度范数裁剪}：

$$
\hat{g} = \begin{cases}
g & \text{if } \|g\| \leq C \\
C \cdot \frac{g}{\|g\|} & \text{if } \|g\| > C
\end{cases}
$$

其中 $C$ 是裁剪阈值（通常设为 1.0），$\|g\|$ 是所有参数梯度拼接后的 L2 范数。

\begin{lstlisting}[caption={梯度裁剪的标准用法}]
loss.backward()
# Clip gradients before optimizer step
grad_norm = torch.nn.utils.clip_grad_norm_(
    model.parameters(), max_norm=1.0
)
optimizer.step()
optimizer.zero_grad(set_to_none=True)

# Monitor grad_norm for training stability
if grad_norm > 10.0:
    print(f”Warning: large gradient norm {grad_norm:.2f}”)
\end{lstlisting}

\begin{importantbox}{梯度范数是训练健康的晴雨表}
监控梯度范数可以帮助诊断训练问题：
\begin{itemize}
    \item 梯度范数持续增大 $\to$ 可能训练发散，需要降低学习率
    \item 梯度范数突然飙升 $\to$ 可能遇到异常数据或数值不稳定
    \item 梯度范数趋近于零 $\to$ 可能学习率太小或梯度消失
    \item 梯度范数频繁被裁剪 $\to$ 裁剪阈值可能设得太低
\end{itemize}
\end{importantbox}

